\documentclass[aspectratio=169,11pt]{beamer}
% Lecture 02 — Java Syntax Basics
% Course: Introduction to Programming with Java

\input{preamble}

\title{Lecture 02: Java Syntax Basics}
\subtitle{Introduction to Programming with Java}

\begin{document}

%----------------------------------------------------------------------
    \begin{frame}
        \titlepage
    \end{frame}

%----------------------------------------------------------------------
    \begin{frame}{Today's learning goals}
        By the end of this lecture you should be able to:
        \begin{enumerate}
            \item Declare and initialise variables with appropriate primitive types
            \item Write \javacode{if}/\javacode{else} and \javacode{switch} decisions
            \item Use \javacode{for}, \javacode{while}, and \javacode{do}-\javacode{while} loops
            \item Create and index arrays; use basic \javacode{String} operations
            \item Read short programs and predict their output
        \end{enumerate}
    \end{frame}

    \begin{frame}{Agenda}
        \begin{enumerate}
            \item Variables and assignment
            \item Primitive data types (focus on integers and friends)
            \item Conditionals: \javacode{if}/\javacode{else} and relational operators
            \item \javacode{switch}
            \item Loops: \javacode{for}, \javacode{while}, \javacode{do}-\javacode{while}
            \item Arrays
            \item \javacode{String} basics
            \item Putting it together
        \end{enumerate}
    \end{frame}

%======================================================================
    \section{Variables}
%======================================================================

    \begin{frame}{What is a variable?}
        A \textbf{variable} is a named storage location for a value of a given type.
        \vspace{0.6em}
        \begin{itemize}
            \item Has a \textbf{name} (identifier)
            \item Has a \textbf{type} (what kind of data it holds)
            \item Holds a \textbf{value} that can change over time (unless \javacode{final})
        \end{itemize}
    \end{frame}

    \begin{frame}[fragile]{Declaring and initialising}
        \begin{lstlisting}
int score;           // declaration
score = 0;           // assignment

int lives = 3;       // declaration + initialisation

final double PI = 3.14159;  // constant (cannot reassign)
        \end{lstlisting}
        \vspace{0.4em}
        \begin{itemize}
            \item Java is \textbf{statically typed}: the type is fixed at compile time
            \item Local variables must be assigned before use
        \end{itemize}
    \end{frame}

    \begin{frame}{Naming rules and style}
        \begin{itemize}
            \item Letters, digits, \javacode{\_}, \javacode{\$}; cannot start with a digit
            \item Case-sensitive: \javacode{score} $\neq$ \javacode{Score}
            \item Convention: \textbf{camelCase} for variables and methods (\javacode{playerScore})
            \item Class names use \textbf{PascalCase} (\javacode{PlayerScore})
            \item Prefer meaningful names over abbreviations
        \end{itemize}
        \vspace{0.4em}
        \commonmistake{Using reserved words as names (\javacode{int}, \javacode{class}, \javacode{for}, \ldots).}
    \end{frame}

    \begin{frame}[fragile]{Assignment and expressions}
        \begin{lstlisting}
int a = 10;
int b = 3;
int sum = a + b;      // 13
int rem = a % b;      // 1 (remainder)
a += 5;               // same as a = a + 5
a++;                  // increment by 1
        \end{lstlisting}
        \vspace{0.3em}
        Operators you will use constantly: \javacode{+} \javacode{-} \javacode{*} \javacode{/} \javacode{\%}
        and compound assignments \javacode{+=} \javacode{-=} \javacode{*=} \javacode{/=}.
    \end{frame}

    \begin{frame}{Integer division (watch out)}
        \begin{itemize}
            \item Dividing two \javacode{int} values truncates toward zero
            \item \javacode{7 / 2} is \javacode{3}, not \javacode{3.5}
            \item Use a floating type if you need a fractional result
        \end{itemize}
        \vspace{0.4em}
        \trythis{Predict \javacode{5 / 2}, \javacode{5 \% 2}, and \javacode{5.0 / 2}. Then verify in a tiny program.}
    \end{frame}

%======================================================================
    \section{Primitive types}
%======================================================================

    \begin{frame}{Primitive types overview}
        Java's eight primitive types:
        \vspace{0.4em}
        \begin{center}
            \begin{tabular}{@{}lll@{}}
                \toprule
                \textbf{Category} & \textbf{Types} & \textbf{Examples} \\
                \midrule
                Integer & \javacode{byte}, \javacode{short}, \javacode{int}, \javacode{long} & \javacode{42}, \javacode{100L} \\
                Floating & \javacode{float}, \javacode{double} & \javacode{3.14}, \javacode{2.5f} \\
                Character & \javacode{char} & \javacode{'A'} \\
                Boolean & \javacode{boolean} & \javacode{true}, \javacode{false} \\
                \bottomrule
            \end{tabular}
        \end{center}
        \vspace{0.5em}
        Default choices for beginners: \javacode{int}, \javacode{double}, \javacode{boolean}, \javacode{char}.
    \end{frame}

    \begin{frame}{Integers in focus}
        \begin{itemize}
            \item \javacode{int} --- 32-bit signed integer; everyday counting and indices
            \item \javacode{long} --- 64-bit; use for larger ranges (suffix \javacode{L})
            \item \javacode{byte} / \javacode{short} --- smaller ranges; less common early on
        \end{itemize}
        \vspace{0.4em}
        \begin{block}{Ranges (approximate mental model)}
            \javacode{int} roughly $\pm 2$ billion;\quad
            \javacode{long} far larger.\\
            Overflow wraps around---do not rely on silent wrap in logic.
        \end{block}
    \end{frame}

    \begin{frame}[fragile]{Other primitives you will meet}
        \begin{lstlisting}
double price = 19.99;
boolean isReady = true;
char grade = 'B';

// float is less common in coursework; prefer double
float ratio = 0.5f;
        \end{lstlisting}
        \vspace{0.3em}
        \begin{itemize}
            \item \javacode{boolean} is only \javacode{true} or \javacode{false} (not 0/1)
            \item \javacode{char} holds a single UTF-16 code unit; use single quotes
            \item \javacode{String} is \emph{not} a primitive (next sections)
        \end{itemize}
    \end{frame}

    \begin{frame}[fragile]{Type casting (gentle intro)}
        \begin{lstlisting}
int whole = 5;
double real = whole;      // widening: OK, automatic

double x = 3.9;
int truncated = (int) x;  // narrowing: 3 (fraction dropped)
        \end{lstlisting}
        \vspace{0.3em}
        \commonmistake{Expecting \javacode{(int) 3.9} to round to 4. It truncates toward zero.}
    \end{frame}

%--- Quiz checkpoint 1 (after Variables + Primitive types) ---
    \begin{frame}{Quiz 1 --- Types and division}
        \quizbox{What does \javacode{System.out.println(7 / 2);} print?}
        \begin{enumerate}
            \item \javacode{3.5}
            \item \javacode{3}
            \item \javacode{4}
            \item A compile error
        \end{enumerate}
        \quizpause
    \end{frame}

    \begin{frame}{Answer --- Quiz 1}
        \answerbox{\textbf{2.} \javacode{3} --- integer division truncates toward zero.}
        \vspace{0.4em}
        Use a floating type (e.g.\ \javacode{7.0 / 2} or \javacode{(double) 7 / 2}) for a fractional result.
    \end{frame}

%======================================================================
    \section{Conditionals}
%======================================================================

    \begin{frame}{Making decisions}
        Programs branch based on conditions:
        \begin{itemize}
            \item Relational operators: \javacode{==} \javacode{!=} \javacode{<} \javacode{>} \javacode{<=} \javacode{>=}
            \item Logical operators: \javacode{\&\&} (and), \javacode{||} (or), \javacode{!} (not)
            \item Result of a comparison is a \javacode{boolean}
        \end{itemize}
    \end{frame}

    \begin{frame}[fragile]{if and else}
        \begin{lstlisting}
int temperature = 18;

if (temperature < 0) {
    System.out.println("Freezing");
} else if (temperature < 15) {
    System.out.println("Cold");
} else {
    System.out.println("Mild or warm");
}
        \end{lstlisting}
    \end{frame}

    \begin{frame}{Braces and structure}
        \begin{itemize}
            \item Always use braces \javacode{\{} \javacode{\}} even for one statement---avoids bugs
            \item Conditions go in parentheses: \javacode{if (score >= 50)}
            \item \javacode{else if} chains test alternatives in order
            \item First matching branch runs; then the chain ends
        \end{itemize}
        \vspace{0.4em}
        \commonmistake{Writing \javacode{if (x = 5)} instead of \javacode{if (x == 5)}.
        Assignment vs equality---Java often catches this because types must be boolean.}
    \end{frame}

    \begin{frame}[fragile]{Logical operators}
        \begin{lstlisting}
int age = 20;
boolean hasId = true;

if (age >= 18 && hasId) {
    System.out.println("Entry allowed");
}

if (age < 12 || age > 65) {
    System.out.println("Discount eligible");
}

if (!hasId) {
    System.out.println("Please show ID");
}
        \end{lstlisting}
    \end{frame}

    \begin{frame}{Short-circuit evaluation}
        \begin{itemize}
            \item \javacode{\&\&} stops if the left side is \javacode{false}
            \item \javacode{||} stops if the left side is \javacode{true}
            \item Useful and important when later conditions would be invalid
        \end{itemize}
        \vspace{0.4em}
        Example idea: check array length before indexing (arrays come soon).
    \end{frame}

%======================================================================
    \section{switch}
%======================================================================

    \begin{frame}[fragile]{switch: multi-way branch}
        \begin{lstlisting}
int day = 3;
switch (day) {
    case 1:
        System.out.println("Monday");
        break;
    case 2:
        System.out.println("Tuesday");
        break;
    case 3:
        System.out.println("Wednesday");
        break;
    default:
        System.out.println("Other day");
}
        \end{lstlisting}
    \end{frame}

    \begin{frame}{switch rules of thumb}
        \begin{itemize}
            \item Classic \javacode{switch} works well with \javacode{int}, \javacode{char}, \javacode{String}, enums
            \item Include \javacode{break} unless you intentionally want fall-through
            \item \javacode{default} handles unmatched values
        \end{itemize}
        \vspace{0.4em}
        \commonmistake{Forgetting \javacode{break} and ``falling through'' into the next case.}
    \end{frame}

    \begin{frame}[fragile]{Modern switch (optional awareness)}
        Java also supports a clearer arrow form (Java~14+):
        \begin{lstlisting}
int day = 3;
String name = switch (day) {
    case 1 -> "Monday";
    case 2 -> "Tuesday";
    case 3 -> "Wednesday";
    default -> "Other day";
};
System.out.println(name);
        \end{lstlisting}
        \vspace{0.3em}
        Use what your course JDK supports; classic \javacode{switch} is still essential.
    \end{frame}

%--- Quiz checkpoint 2 (after Conditionals + switch) ---
    \begin{frame}[fragile]{Quiz 2 --- What does this print?}
        \begin{lstlisting}
int x = 10;
if (x > 15) {
    System.out.println("A");
} else if (x > 5) {
    System.out.println("B");
} else {
    System.out.println("C");
}
        \end{lstlisting}
        \quizbox{Output?}
        \begin{enumerate}
            \item \javacode{A}
            \item \javacode{B}
            \item \javacode{C}
            \item Nothing (no branch runs)
        \end{enumerate}
        \quizpause
    \end{frame}

    \begin{frame}{Answer --- Quiz 2}
        \answerbox{\textbf{2.} \javacode{B} --- \javacode{x > 15} is false; \javacode{x > 5} is true, so that branch runs and the chain stops.}
    \end{frame}

    \begin{frame}{Quiz 2b --- switch (true/false)}
        \quizbox{True or false?\\[0.4em]
        In a classic \javacode{switch}, forgetting \javacode{break} can cause execution to
        continue into the next \javacode{case} (fall-through).}
        \begin{enumerate}
            \item True
            \item False
        \end{enumerate}
        \quizpause
    \end{frame}

    \begin{frame}{Answer --- Quiz 2b}
        \answerbox{\textbf{True.} Always add \javacode{break} unless you intentionally want fall-through.}
    \end{frame}

%======================================================================
    \section{Loops}
%======================================================================

    \begin{frame}{Why loops?}
        Repeat work without copying code:
        \begin{itemize}
            \item Process every element of a collection
            \item Repeat until a condition becomes false
            \item Count from a start value to an end value
        \end{itemize}
        \vspace{0.4em}
        Three core forms: \javacode{for}, \javacode{while}, \javacode{do}-\javacode{while}.
    \end{frame}

    \begin{frame}[fragile]{for loop}
        \begin{lstlisting}
for (int i = 0; i < 5; i++) {
    System.out.println(i);
}
// prints 0 1 2 3 4  (each on its own line)
        \end{lstlisting}
        \vspace{0.3em}
        Anatomy: \textbf{initialisation}; \textbf{condition}; \textbf{update}.
        \begin{itemize}
            \item Prefer starting at \javacode{0} when working with arrays later
            \item Condition checked \emph{before} each iteration
        \end{itemize}
    \end{frame}

    \begin{frame}[fragile]{while loop}
        \begin{lstlisting}
int n = 3;
while (n > 0) {
    System.out.println(n);
    n--;
}
System.out.println("Go!");
        \end{lstlisting}
        \vspace{0.3em}
        Use \javacode{while} when the number of iterations is not obvious in advance.
        \vspace{0.3em}
        \commonmistake{Forgetting to update the loop variable $\rightarrow$ infinite loop.}
    \end{frame}

    \begin{frame}[fragile]{do-while loop}
        \begin{lstlisting}
int guess;
do {
    guess = readGuess();   // pretend this reads input
} while (guess < 1 || guess > 10);
        \end{lstlisting}
        \vspace{0.3em}
        \begin{itemize}
            \item Body runs \textbf{at least once}
            \item Condition checked \emph{after} the body
            \item Handy for ``ask until valid'' patterns
        \end{itemize}
    \end{frame}

    \begin{frame}[fragile]{Choosing a loop}
        \begin{tabular}{@{}lp{0.55\textwidth}@{}}
            \toprule
            \textbf{Prefer} & \textbf{When} \\
            \midrule
            \javacode{for} & Counting; known range; array indices \\
            \javacode{while} & Repeat while a condition holds \\
            \javacode{do}-\javacode{while} & Must run once; then maybe again \\
            \bottomrule
        \end{tabular}
        \vspace{0.6em}
        \trythis{Rewrite a \javacode{for} that prints 1..10 as a \javacode{while}. Same output?}
    \end{frame}

    \begin{frame}[fragile]{break and continue (brief)}
        \begin{lstlisting}
for (int i = 1; i <= 10; i++) {
    if (i == 3) {
        continue;   // skip the rest of this iteration
    }
    if (i == 8) {
        break;      // leave the loop entirely
    }
    System.out.println(i);
}
        \end{lstlisting}
        Use sparingly; clear loop conditions are often easier to read.
    \end{frame}

%--- Quiz checkpoint 3 (after Loops) ---
    \begin{frame}[fragile]{Quiz 3 --- Loop output}
        \begin{lstlisting}
int n = 3;
while (n > 0) {
    System.out.print(n);
    n--;
}
        \end{lstlisting}
        \quizbox{What is printed?}
        \begin{enumerate}
            \item \javacode{321}
            \item \javacode{333}
            \item \javacode{012}
            \item Nothing (infinite loop, never prints)
        \end{enumerate}
        \quizpause
    \end{frame}

    \begin{frame}{Answer --- Quiz 3}
        \answerbox{\textbf{1.} \javacode{321} --- prints while \javacode{n} is 3, then 2, then 1; then \javacode{n} becomes 0 and the loop ends.}
    \end{frame}

%======================================================================
    \section{Arrays}
%======================================================================

    \begin{frame}{What is an array?}
        An \textbf{array} stores a fixed number of values of the \emph{same} type,
        accessed by an integer index.
        \vspace{0.5em}
        \begin{itemize}
            \item Length is fixed after creation
            \item Indices start at \javacode{0}
            \item Last index is \javacode{length - 1}
        \end{itemize}
    \end{frame}

    \begin{frame}[fragile]{Creating and using arrays}
        \begin{lstlisting}
int[] scores = new int[3];  // {0, 0, 0}
scores[0] = 90;
scores[1] = 85;
scores[2] = 92;

int[] primes = {2, 3, 5, 7};
System.out.println(primes.length);  // 4
System.out.println(primes[0]);      // 2
        \end{lstlisting}
    \end{frame}

    \begin{frame}[fragile]{Traversing an array}
        \begin{lstlisting}
int[] values = {4, 1, 9, 2};
int sum = 0;

for (int i = 0; i < values.length; i++) {
    sum += values[i];
}
System.out.println(sum);  // 16

// enhanced for (read each element):
for (int v : values) {
    System.out.println(v);
}
        \end{lstlisting}
    \end{frame}

    \begin{frame}{Bounds and defaults}
        \begin{itemize}
            \item Accessing \javacode{values[values.length]} throws
            \javacode{ArrayIndexOutOfBoundsException}
            \item Numeric arrays default to \javacode{0}; \javacode{boolean} to \javacode{false};
            reference types to \javacode{null}
        \end{itemize}
        \vspace{0.4em}
        \commonmistake{Using \javacode{<= values.length} in a \javacode{for} loop header.
        Use \javacode{< values.length}.}
    \end{frame}

    \begin{frame}[fragile]{Arrays of other types}
        \begin{lstlisting}
double[] temps = {18.5, 21.0, 19.2};
boolean[] flags = {true, false, true};
char[] letters = {'J', 'a', 'v', 'a'};
String[] names = {"Ada", "Alan", "Grace"};
        \end{lstlisting}
        \vspace{0.3em}
        Multidimensional arrays exist (\javacode{int[][]}), but master 1D arrays first.
    \end{frame}

%--- Quiz checkpoint 4 (after Arrays) ---
    \begin{frame}[fragile]{Quiz 4 --- Arrays}
        \begin{lstlisting}
int[] a = {10, 20, 30};
System.out.println(a.length);
System.out.println(a[0]);
        \end{lstlisting}
        \quizbox{What is printed (two lines)?}
        \begin{enumerate}
            \item \javacode{2} then \javacode{10}
            \item \javacode{3} then \javacode{10}
            \item \javacode{3} then \javacode{20}
            \item \javacode{3} then a compile error
        \end{enumerate}
        \quizpause
    \end{frame}

    \begin{frame}{Answer --- Quiz 4}
        \answerbox{\textbf{2.} \javacode{3} then \javacode{10} --- length is 3; index \javacode{0} is the first element.}
        \vspace{0.3em}
        Last valid index is \javacode{a.length - 1} (here: \javacode{2}).
    \end{frame}

%======================================================================
    \section{String basics}
%======================================================================

    \begin{frame}{String is a class}
        \begin{itemize}
            \item \javacode{String} holds text: sequences of characters
            \item Not a primitive---it is a reference type (objects)
            \item Literals use double quotes: \javacode{"Hello"}
            \item Strings are \textbf{immutable}: operations create new strings
        \end{itemize}
    \end{frame}

    \begin{frame}[fragile]{Creating and concatenating}
        \begin{lstlisting}
String first = "Hello";
String second = "Java";
String msg = first + ", " + second + "!";
System.out.println(msg);           // Hello, Java!

int n = 7;
System.out.println("Count: " + n); // Count: 7
        \end{lstlisting}
    \end{frame}

    \begin{frame}[fragile]{Useful String methods}
        \begin{lstlisting}
String s = "Java";
System.out.println(s.length());        // 4
System.out.println(s.charAt(0));       // J
System.out.println(s.toUpperCase());   // JAVA
System.out.println(s.substring(1, 3)); // av

String t = "java";
System.out.println(s.equals(t));           // false
System.out.println(s.equalsIgnoreCase(t)); // true
        \end{lstlisting}
    \end{frame}

    \begin{frame}[fragile]{Comparing strings correctly}
        \begin{lstlisting}
String a = "hi";
String b = "hi";
String c = new String("hi");

// Prefer equals for content:
System.out.println(a.equals(b));  // true
System.out.println(a.equals(c));  // true

// == compares references (identity), not always content
        \end{lstlisting}
        \vspace{0.3em}
        \commonmistake{Using \javacode{==} to compare \javacode{String} contents.
        Use \javacode{equals} (or \javacode{equalsIgnoreCase}).}
    \end{frame}

    \begin{frame}{char vs String}
        \begin{tabular}{@{}lll@{}}
            \toprule
            & \textbf{\javacode{char}} & \textbf{\javacode{String}} \\
            \midrule
            Quotes & single \javacode{'A'} & double \javacode{"A"} \\
            Meaning & one character & text (any length) \\
            Kind & primitive & object \\
            \bottomrule
        \end{tabular}
    \end{frame}

%--- Quiz checkpoint 5 (after String basics) ---
    \begin{frame}{Quiz 5 --- Comparing strings}
        \quizbox{Which is the recommended way to compare the \emph{text content} of two \javacode{String} values \javacode{s} and \javacode{t}?}
        \begin{enumerate}
            \item \javacode{s == t}
            \item \javacode{s.equals(t)}
            \item \javacode{s = t}
            \item \javacode{s.compare(t)}
        \end{enumerate}
        \quizpause
    \end{frame}

    \begin{frame}{Answer --- Quiz 5}
        \answerbox{\textbf{2.} \javacode{s.equals(t)} (or \javacode{equalsIgnoreCase} when case should not matter).}
        \vspace{0.4em}
        \javacode{==} compares references (identity), which is not reliable for string \emph{content}.
    \end{frame}

%======================================================================
    \section{Putting it together}
%======================================================================

    \begin{frame}[fragile]{Mini example: average of scores}
        \begin{lstlisting}
int[] scores = {80, 90, 70, 100};
int sum = 0;
for (int i = 0; i < scores.length; i++) {
    sum += scores[i];
}
double average = (double) sum / scores.length;
System.out.println("Average: " + average);

if (average >= 90) {
    System.out.println("Excellent");
} else if (average >= 70) {
    System.out.println("Good");
} else {
    System.out.println("Keep practising");
}
        \end{lstlisting}
    \end{frame}

    \begin{frame}[fragile]{Mini example: count vowels}
        \begin{lstlisting}
String word = "Education";
int vowels = 0;
String lower = word.toLowerCase();

for (int i = 0; i < lower.length(); i++) {
    char c = lower.charAt(i);
    if (c == 'a' || c == 'e' || c == 'i'
            || c == 'o' || c == 'u') {
        vowels++;
    }
}
System.out.println("Vowels: " + vowels);
        \end{lstlisting}
    \end{frame}

    \begin{frame}{Debugging checklist}
        When something fails:
        \begin{enumerate}
            \item Read the compiler error: line number + message
            \item Check types and missing semicolons / braces
            \item Trace loop bounds and array indices
            \item Print intermediate values (\javacode{System.out.println})
            \item Change one thing; recompile; rerun
        \end{enumerate}
    \end{frame}

    \begin{frame}{Key takeaways}
        \begin{enumerate}
            \item Variables need a type, a name, and a defined value before use
            \item Prefer \javacode{int}, \javacode{double}, \javacode{boolean}, \javacode{char} early on
            \item \javacode{if}/\javacode{else} and \javacode{switch} express decisions
            \item \javacode{for}/\javacode{while}/\javacode{do}-\javacode{while} express repetition
            \item Arrays: fixed length, 0-based indices
            \item Strings: use methods and \javacode{equals}, not \javacode{==} for content
        \end{enumerate}
    \end{frame}

    \begin{frame}{Practice before next time}
        \trythis{Write a program that reads (or hard-codes) five integers into an array,
            prints the maximum, and prints whether the maximum is even or odd.}
        \vspace{0.6em}
        \begin{block}{Coming next}
            Methods, parameters, and structuring larger programs---
            building on today's control flow and data basics.
        \end{block}
    \end{frame}

    \begin{frame}{Questions?}
        \begin{center}
            {\Large Which construct do you want to see one more example of?}
        \end{center}
    \end{frame}

\end{document}
